% Chapter 8 -- Working with the Accessories: a worked example (the
% What's New in System 7 habit of ending a topic with the reader's own
% task).  Sources (Law 1): os/apps/calcda (the Calculator; its Edit menu
% dimmed, so a result is typed, not pasted); os/apps/itext/itext.c (typing,
% Save into a folder through Standard File: a double-click on a folder in
% the list opens it, kmain.c sf_run); os/flair/tbxgate.c (one disk
% application at a time: the Calculator is closed before InitechText
% opens).
% GATES: test-flair-calcda test-flair-itext test-stdfile test-flair-app-launch test-pubs-flair-screens
\Gates{test-flair-calcda, test-flair-itext, test-stdfile, test-flair-app-launch,
  test-pubs-flair-screens}
\chapter{Working with the Accessories}

This chapter puts the accessories to work on a small task: working out a
figure with the Calculator and sending it round in a memo written with
InitechText, saved in the \file{APPS} folder.

\section{Working out the figure}
\begin{TaskSteps}
\Step Choose \menuitem{Calculator} from the Apple menu.
\Step Type \type{1200*3=}.
\Result{The display shows \texttt{3600}.}
\end{TaskSteps}

\FlairScreen[scale=4, label=fig:calcmemo]{a_calcmemo}{1200 * 3 = 3600}

\begin{TaskSteps}[resume]
\Step Write the figure down, then click the Calculator's close box.
\Result{The Calculator's \menuitem{Copy} command is dimmed in this
version, so the figure is carried to the memo by hand.  The Calculator must
be closed before InitechText can open: only one application from the disk
is open at a time.}
\end{TaskSteps}

\section{Writing the memo}
\begin{TaskSteps}
\Step Open the \file{APPS} folder and double-click \file{TEXT.EXE}.
\Step Type the memo, with the figure in it.
\Step Press \key{Ctrl-S}.
\Result{The Save dialog box asks for a name.  It lists the top level of the
data disk.}
\Step Double-click \file{APPS} in the list.
\Result{The list now shows the contents of the \file{APPS} folder, and the
box at the top names it.  The document will be saved there.}
\end{TaskSteps}

\FlairScreen[label=fig:sfapps]{a_sfapps}{The Save dialog box, opened at the
APPS folder}

\begin{TaskSteps}[resume]
\Step Erase the name in the name box with \key{Backspace}, type a name such
as \type{memo.txt}, and press \key{Enter}.
\Result{The memo is saved in the \file{APPS} folder, and the window takes
its name.  Its icon appears in the \texttt{APPS} window the next time the
window is opened.}
\Step Press \key{Ctrl-Q} to quit InitechText.
\end{TaskSteps}

\section{Where to go from here}
The \doctitle{FLAIR User's Guide} describes the desktop, the Finder and the
Trash, and Chapter~14 of it tells you how to copy your memo to another disk
at the InitechDOS prompt.
